% chapters/08/section-02.tex
\section{Persepsi Jarak dan Arah}
\label{sec:persepsi}

Langkah pertama pada susunan keputusan adalah persepsi. Pada bab ini, yang
diamati hanyalah posisi pemain, dan yang dihitung dari keterangan itu ada dua:
jarak dan arah.

\subsection{Vektor Menuju Pemain}

Jarak dan arah diperoleh dari satu perhitungan yang sama, yaitu selisih kedua
posisi.

\begin{kodekeluaran}
to_player = player.position - self.position
self.distance = to_player.length()
\end{kodekeluaran}

Baris pertama menghasilkan vektor yang menunjuk dari NPC ke pemain. Perhatikan
urutan pengurangannya: yang dikurangi adalah posisi NPC, sehingga vektornya
menunjuk \textbf{dari NPC menuju pemain}. Apabila urutannya dibalik, vektornya
menunjuk arah sebaliknya dan NPC justru menjauh. Kekeliruan itu sering terjadi
karena kedua urutan sama-sama masuk akal ketika ditulis.

Baris kedua menghitung panjang vektor itu, yaitu jaraknya dalam piksel. Nilai
itulah yang dipakai untuk memutuskan, sedangkan vektornya dipakai untuk
bergerak. Perbedaan peran keduanya diperlihatkan pada
Gambar~\ref{fig:persepsi-npc}.

\begin{figure}[htbp]
\centering
% chapters/08/figures/persepsi-npc.tikz
% Persepsi NPC: jarak ke pemain menentukan keputusan, arah menuju pemain
% menentukan kecepatan, dan besarnya kecepatan tetap berapa pun jaraknya.
\begin{tikzpicture}[
  npc/.style={rectangle, draw=ubRed, thick, fill=ubRed!25, inner sep=4pt},
  pemain/.style={rectangle, draw=ubPrimary, thick, fill=ubPrimary!25, inner sep=4pt},
  lingkar/.style={draw=ubLine, thick, dashed},
  lingkarkecil/.style={draw=ubGreen, thick},
  garisjarak/.style={draw=ubWarn, thick, dashed},
  panah/.style={-{Stealth[length=2.4mm]}, very thick, draw=ubAccent},
  label/.style={font=\scriptsize, align=left}
]

% NPC di titik asal, dengan lingkaran persepsi dan lingkaran jarak henti.
\draw[lingkar] (0,0) circle (3.0);
\draw[lingkarkecil] (0,0) circle (0.55);
\node[npc] (npc) at (0,0) {NPC};

% Pemain berada di dalam jari-jari persepsi.
\node[pemain] (player) at (2.3,1.8) {P};

% Garis jarak dan garis arah.
\draw[garisjarak] (npc.center) -- (player.center);
\node[label, ubWarn, anchor=south east] at (1.35,1.15) {$d$};

\draw[panah] (0.3,0.24) -- (1.15,0.9);
\node[label, ubAccent, anchor=west] at (1.05,0.72)
  {kecepatan $= \dfrac{\text{arah}}{|d|} \times v$};

% Keterangan lingkaran.
\node[label, ubLine, anchor=west] at (2.15,2.15) {jari-jari persepsi};
\draw[ubLine, thin] (2.1,2.1) -- (2.0,1.35);

\node[label, ubGreen, anchor=north west] at (0.5,-0.75) {jarak henti};
\draw[ubGreen, thin] (0.55,-0.7) -- (0.38,-0.42);

% Keterangan perilaku.
\node[label, anchor=north, text width=104mm] at (0,-3.5)
  {Apabila pemain berada di dalam jari-jari persepsi dan di luar jarak henti,
   NPC bergerak menuju pemain dengan kecepatan tetap. Perhatikan bahwa panjang
   panah kecepatan tidak bergantung pada jarak $d$.};
\end{tikzpicture}

\caption{Persepsi NPC: jarak untuk memutuskan, arah untuk bergerak}
\label{fig:persepsi-npc}
\end{figure}

\subsection{Jari-jari Persepsi sebagai Batas Keputusan}

NPC tidak seharusnya mengetahui keberadaan pemain di mana pun. Jangkauan
penglihatannya dibatasi oleh sebuah nilai yang disebut jari-jari persepsi.
Nilai itu dibandingkan dengan jarak yang baru saja dihitung.

\begin{kodekeluaran}
if self.distance <= self.radius:
    # pemain terlihat: keputusan berikutnya adalah mengejar
else:
    # pemain tidak terlihat: berdiam
\end{kodekeluaran}

Nilai jari-jari itu bukan urusan teknis, melainkan urusan perancangan. Jari-jari
yang terlalu besar membuat NPC terasa mengetahui segala hal; jari-jari yang
terlalu kecil membuat NPC mudah dihindari sehingga permainan kehilangan
tantangan. Pada program bab ini, nilainya dapat diubah melalui baris perintah
agar pengaruhnya dapat dirasakan langsung.

\begin{tipsbox}
Untuk membandingkan jarak dengan jari-jari, kuadrat jaraknya
sudah mencukupi, sehingga akar pada \code{length()} dapat dihindari. Pada bab
ini \code{length()} tetap dipakai karena jaraknya juga ditampilkan pada layar
dan dipakai untuk keperluan lain. Kebiasaan memilih perhitungan yang paling
murah---sebagaimana disebutkan pada Bagian~5.2---baru menjadi penting ketika
jumlah NPC bertambah banyak, dan hal itu dibahas pada Bab~XI.
\end{tipsbox}

\subsection{Menggambar Persepsi agar Keputusannya Terlihat}

Perilaku NPC sulit diperiksa apabila hanya dilihat dari geraknya. Karena itu,
program bab ini menggambar tiga hal yang memperlihatkan keputusannya.

\begin{itemize}
  \item \textbf{Lingkaran jari-jari persepsi}, sehingga batas jangkauannya
        tampak.
  \item \textbf{Lingkaran jarak henti}, yaitu batas kapan NPC berhenti,
        sebagaimana dibahas pada Bagian~\ref{sec:berhenti}.
  \item \textbf{Garis arah menuju pemain}, yaitu garis tipis dari setiap NPC ke
        pemain. Garis ini bukan bagian dari permainan yang sesungguhnya; ia
        dipakai untuk memeriksa apakah arah yang dihitung memang benar.
  \item \textbf{Warna NPC}, yang berubah menurut keputusannya: satu warna
        ketika mengejar, warna lain ketika berdiam.
\end{itemize}

Cara menggambar keadaan internal semacam ini merupakan kebiasaan yang sangat
berguna ketika mengembangkan kecerdasan buatan. Tanpa gambar itu, Anda hanya
melihat NPC yang bergerak atau tidak bergerak, dan tidak dapat membedakan
``berdiam karena memutuskan berdiam'' dari ``berdiam karena kesalahan
perhitungan''.
